\documentclass[aps,prd,twocolumn,superscriptaddress,nofootinbib]{revtex4-2}

\usepackage{amsmath,amssymb,amsfonts}
\usepackage{graphicx}
\usepackage{hyperref}
\usepackage{bm}
\usepackage{booktabs}
\usepackage{amsthm}

\newtheorem{theorem}{Theorem}
\newtheorem{proposition}[theorem]{Proposition}
\newtheorem{conjecture}[theorem]{Conjecture}

\newcommand{\Spec}{\mathrm{Spec}}
\newcommand{\Tr}{\mathrm{Tr}}
\newcommand{\Br}{\mathrm{Br}}
\newcommand{\Gal}{\mathrm{Gal}}
\newcommand{\Zp}{\mathbb{Z}[1/p]}
\newcommand{\Qp}{\mathbb{Q}_p}
\newcommand{\Zhat}{\hat{\mathbb{Z}}}
\newcommand{\Fp}{\mathbb{F}_p}

\begin{document}

\title{From Bost-Connes to Warp Drive:\\
  A Theoretical Chain via Noncommutative Geometry\\
  and Arithmetic Vacuum Engineering}

\author{Wright Brothers}
\affiliation{Independent Research}
\date{\today}

\begin{abstract}
We present a complete theoretical chain connecting the Bost-Connes
quantum statistical mechanical system to the possibility of
Alcubierre warp drive, via Connes' noncommutative geometry (NCG)
program. The chain consists of ten steps, of which we classify
seven as established, two as argued (relying on published but
not yet experimentally confirmed theorems), and one as open
(the scaling problem). We give a precise formulation of each step,
identify the mathematical theorem or conjecture supporting it,
and define the open problems that must be resolved for the chain
to be complete. The key open problem --- bridging the energy gap
of $\sim 10^{68}$ orders of magnitude between laboratory-scale
vacuum energy modifications and macroscopic spacetime curvature ---
is formulated as five distinct sub-problems, each corresponding to
a candidate amplification mechanism. This paper serves as a roadmap
for the arithmetic vacuum engineering research program introduced
in the companion paper~\cite{paper_A}.
\end{abstract}

\maketitle

% ============================================================================
\section{Introduction}
\label{sec:intro}
% ============================================================================

In a companion paper~\cite{paper_A}, we showed that ``prime channel
muting'' in the Bost-Connes (BC) system flips the sign of
zeta-regularized vacuum energy, potentially violating the weak
energy condition (WEC) required by the Alcubierre warp
metric~\cite{Alcubierre1994}. That paper focused on the algebraic
mechanism and a concrete experimental proposal.

The present paper addresses a different question: \emph{what is the
complete theoretical chain from the BC system to actual spacetime
modification, and where are the gaps?}

We identify a chain of ten logical steps (Table~\ref{tab:chain})
and classify each according to its epistemic status:
\textbf{Established} (proven mathematically or experimentally),
\textbf{Argued} (supported by published theorems whose physical
applicability is conjectured), or \textbf{Open} (no satisfactory
argument exists). We give precise references for each step and
formulate the open problems as research targets.

\begin{table*}[t]
\centering
\caption{The ten-step theoretical chain from the Bost-Connes system
to warp drive. Status: E = Established, A = Argued, O = Open.}
\label{tab:chain}
\begin{tabular}{@{}clllp{5.5cm}@{}}
\toprule
Step & Name & Status & Reference & Content \\
\midrule
1 & NCG Standard Model & E & \cite{Connes1996,CCM2007} & $M^4 \times F$: spacetime $=$ 4-manifold $\times$ finite NC space \\
2 & BC as arithmetic kernel & A & \cite{ConnesMarcolli2006} & BC $C^*$-algebra encodes arithmetic of $F$ \\
3 & Phase transition & E & \cite{BostConnes1995} & $\beta = 1$: spontaneous breaking of $\Gal(\mathbb{Q}^{\mathrm{ab}}/\mathbb{Q})$ \\
4 & Prime muting & E & \cite{paper_A} & $P_{\neg p}$: Hilbert space projection, $Z_{\neg p} = \zeta \cdot (1-p^{-\beta})$ \\
5 & Partition function identity & E & Euler product & $Z_{\neg p}(\beta) = \zeta(\beta)(1-p^{-\beta})$, algebraic identity \\
6 & Vacuum energy sign flip & E & \cite{paper_A} & $\zeta_{\neg p}(-3) = (1-p^3)/120 < 0$ for any $p \geq 2$ \\
7 & WEC violation & E & Standard QFT & Negative $\zeta$-regularized energy $\Rightarrow$ $\rho < 0$ locally \\
8 & Warp metric allowed & E & \cite{Alcubierre1994} & $\rho < 0$ is sufficient for Alcubierre $ds^2$ \\
9 & NCG $\leftrightarrow$ bulk spacetime & A & \cite{ConnesMarcolli2006,Maldacena1998} & BC operations on $F$ $\leftrightarrow$ curvature on $M^4$ \\
10 & Scaling & O & --- & $10^{-22}\,\mathrm{J} \to 10^{47}\,\mathrm{J}$ \\
\bottomrule
\end{tabular}
\end{table*}

% ============================================================================
\section{Established Steps (1, 3--8)}
\label{sec:established}
% ============================================================================

\subsection{Step 1: NCG Standard Model}

Connes' noncommutative geometry (NCG) approach to particle
physics~\cite{Connes1996,CCM2007} describes spacetime as a product
\begin{equation}
  \mathcal{M} = M^4 \times F,
\end{equation}
where $M^4$ is a compact Riemannian spin 4-manifold and
$F$ is a finite-dimensional noncommutative space described by the
spectral triple $(\mathcal{A}_F, \mathcal{H}_F, D_F)$ with
$\mathcal{A}_F = \mathbb{C} \oplus \mathbb{H} \oplus M_3(\mathbb{C})$.

The spectral action principle~\cite{CCM2007}
$S = \Tr(f(D/\Lambda))$
reproduces the full Standard Model Lagrangian coupled to
Einstein-Hilbert gravity, including the Higgs field as the
inner fluctuation of the metric on $F$.
This is an established mathematical result.

\subsection{Steps 3--6: From BC system to sign flip}

These steps form the core of the companion paper~\cite{paper_A}
and are summarized here.

\textbf{Step 3.} The BC system~\cite{BostConnes1995} on
$\ell^2(\mathbb{N}^*)$ with $H|n\rangle = \log(n)|n\rangle$
has partition function $Z(\beta) = \zeta(\beta)$ and a phase
transition at $\beta = 1$ with symmetry group
$\Gal(\mathbb{Q}^{\mathrm{ab}}/\mathbb{Q})$.

\textbf{Step 4.} The projection $P_{\neg p} = \sum_{p \nmid n} |n\rangle\langle n|$
defines prime channel muting.

\textbf{Step 5.} The identity
$Z_{\neg p}(\beta) = \zeta(\beta)(1-p^{-\beta})$
follows from the Euler product and is exact.

\textbf{Step 6.} Evaluating at $s = -3$:
$\zeta_{\neg p}(-3) = (1-p^3)/120 < 0$ for all primes $p \geq 2$.

\subsection{Steps 7--8: From negative energy to warp metric}

\textbf{Step 7.} If the zeta-regularized vacuum energy density is
negative, the WEC ($T_{\mu\nu} u^\mu u^\nu \geq 0$) is violated.
This is the standard interpretation of the Casimir
effect~\cite{Casimir1948,Bordag2009}: boundary conditions modify
vacuum energy, and zeta regularization yields the physical
(finite, measurable) result.

\textbf{Step 8.} The Alcubierre metric~\cite{Alcubierre1994} requires
$\rho < 0$ in the bubble wall. Step~7 provides this.
This step is tautological: if negative energy exists, the metric
is a valid solution of Einstein's equations.

% ============================================================================
\section{Argued Steps (2, 9)}
\label{sec:argued}
% ============================================================================

These steps rely on established mathematical results whose
\emph{physical} interpretation is conjectural.

\subsection{Step 2: BC system as arithmetic kernel of $F$}

\begin{proposition}[Connes-Marcolli~\cite{ConnesMarcolli2006}]
The Bost-Connes $C^*$-dynamical system is the noncommutative
space of commensurability classes of $\mathbb{Q}$-lattices
modulo scaling:
\begin{equation}
  \mathcal{A}_{\mathrm{BC}} = C^*(\mathbb{Q}_+^* \backslash \mathbb{A}_f / \Zhat^*),
  \label{eq:bc_algebra}
\end{equation}
where $\mathbb{A}_f$ denotes the finite adeles of $\mathbb{Q}$.
\end{proposition}

This is Theorem~4.1 of Connes-Marcolli~\cite{ConnesMarcolli2006},
established mathematically. The \emph{physical} claim is:

\begin{conjecture}[Arithmetic kernel]
\label{conj:kernel}
The internal space $F$ in Connes' $M^4 \times F$ contains
$\mathcal{A}_{\mathrm{BC}}$ as an arithmetic subalgebra
governing the number-theoretic structure of vacuum states.
Specifically, the KMS states of $\mathcal{A}_{\mathrm{BC}}$
at inverse temperature $\beta$ correspond to vacuum sectors
of quantum fields on $M^4$.
\end{conjecture}

\textbf{Evidence for the conjecture:}

\textit{(i) Structural parallel.}
Both $F$ and $\mathcal{A}_{\mathrm{BC}}$ are
noncommutative $C^*$-algebras with similar structural
properties (finite-dimensional representations,
automorphism groups related to number fields).

\textit{(ii) Spectral action and $\zeta$.}
The spectral action on $M^4 \times F$ produces
$\zeta$-function values~\cite{CCM2007}, and the BC
partition function \emph{is} the $\zeta$ function.

\textit{(iii) Connes-Marcolli program.}
Connes-Marcolli~\cite{ConnesMarcolli2006} explicitly
develop the BC system as an arithmetic model for quantum
statistical mechanics ``over $\Spec(\mathbb{Z})$.''

Beyond these structural arguments, we submit that there is
a deeper, \emph{Einsteinian} argument for why the physical
vacuum should possess BC structure:

\textit{(iv) The unreasonable effectiveness of $\zeta$-regularization.}
Zeta function regularization assigns the value
$\zeta(-1) = -1/12$ to the divergent sum $1+2+3+\cdots$,
and $\zeta(-3) = 1/120$ to $\sum n^3$.
These values, used in Casimir energy calculations, yield
predictions that agree with experiment to within
$1\%$~\cite{Lamoreaux1997}.
If the vacuum were merely ``a collection of field modes''
with no arithmetic structure, there would be no reason for
a number-theoretic regularization to produce physically
correct results. The empirical success of $\zeta$-regularization
is, in Einstein's sense, a ``miracle'' that demands structural
explanation. The simplest explanation: the vacuum \emph{is}
described by $\zeta(s)$, because its underlying structure
is arithmetic.

\textit{(v) The Euler product is not a choice but a consequence.}
The factorization $\zeta(s) = \prod_p (1-p^{-s})^{-1}$
is equivalent to the fundamental theorem of arithmetic
(unique prime factorization). It is not one of many possible
decompositions of $\zeta(s)$; it is the \emph{unique}
multiplicative decomposition. If vacuum energy equals $\zeta(-3)$,
then it \emph{already} has prime structure --- we do not impose
it; we read it.

\textit{(vi) The cosmological constant problem as diagnostic.}
Naive cutoff regularization of vacuum energy yields a value
$\sim 10^{120}$ times the observed cosmological constant ---
the worst prediction in the history of physics.
Zeta regularization, by contrast, yields a finite value
($\zeta(-3) = 1/120$). The fact that the ``wrong'' (cutoff)
regularization diverges while the ``right'' ($\zeta$)
regularization converges suggests that the vacuum is not
well-described as a sum-with-cutoff but \emph{is}
well-described as a $\zeta$-function. And the natural quantum
mechanical system whose partition function is $\zeta$ is
precisely the BC system.

\textit{(vii) Einstein's precedent: geometry as primary.}
Einstein's central insight in general relativity was that
gravity is not a force acting in spacetime but a property of
spacetime geometry itself. By analogy: $\zeta$-regularization
is not a computational trick applied to vacuum energy but a
reflection of the arithmetic-geometric structure of the vacuum
itself. Just as Einstein replaced ``force in flat space'' with
``curvature of space,'' we propose replacing ``regularization
of field modes'' with ``arithmetic structure of the vacuum
over $\Spec(\mathbb{Z})$.''

\textbf{What is missing:} An explicit embedding
$\mathcal{A}_{\mathrm{BC}} \hookrightarrow \mathcal{A}_F$
or a derivation showing that operations on
$\mathcal{A}_{\mathrm{BC}}$ induce specific modifications of
the spectral action on $M^4$. The arguments (iv)--(vii) above
are heuristic --- they explain \emph{why one should expect}
the conjecture to hold, not \emph{that} it holds.
However, we note that Einstein's own arguments for general
relativity were similarly heuristic (the equivalence principle,
Mach's principle) before the mathematical framework was
established. The experimental test proposed in~\cite{paper_A}
is designed to probe exactly this question.

\subsection{Step 9: NCG operations $\leftrightarrow$ bulk curvature}

The claim that modifying the BC partition function corresponds
to modifying the spacetime geometry on $M^4$ requires a bridge
between the algebraic operation (prime muting) and the geometric
consequence (spacetime curvature).

Two bridges have been proposed:

\textbf{Bridge A: Spectral action.}
The spectral action $\Tr(f(D/\Lambda))$ on $M^4 \times F$
produces the Einstein-Hilbert action plus matter terms.
If the BC system modifies $F$, the spectral action changes,
and the equations of motion (including Einstein's equations)
are modified. This is the most direct route but requires
computing how $P_{\neg p}$ affects $D_F$.

\textbf{Bridge B: Holographic.}
The BC system is a 1-dimensional quantum mechanical system.
Via holographic (AdS/CFT-type) correspondence~\cite{Maldacena1998},
a boundary quantum system can encode bulk gravitational physics.
If the BC system serves as a ``boundary theory'' for an
arithmetic bulk, its partition function modification
$\zeta \to \zeta_{\neg p}$ would correspond to a bulk geometry
modification.

Both bridges are conjectural. Neither has been derived from
first principles.

% ============================================================================
\section{The Open Problem: Scaling (Step 10)}
\label{sec:scaling}
% ============================================================================

The most critical gap in the chain is quantitative.

\subsection{Statement of the problem}

The vacuum energy modification measurable in a laboratory
experiment (Section~VI of~\cite{paper_A}) is
$|\Delta E| \sim 10^{-22}\,\mathrm{J}$.
The energy required for an Alcubierre bubble
($R = 50\,\mathrm{m}$, $v = c$, wall thickness $\sim 5\,\mathrm{m}$)
is $\sim 10^{47}\,\mathrm{J}$.
The gap is $\sim 10^{69}$ orders of magnitude.

Even at the Planck scale, the minimum warp bubble
requires $E \sim E_P \sim 10^9\,\mathrm{J}$, a gap
of $\sim 10^{31}$ from laboratory scales.

\subsection{Five candidate mechanisms}

We formulate the scaling problem as five distinct sub-problems,
each corresponding to a candidate amplification mechanism.

\textbf{Mechanism I: Coherent amplification.}
If $N$ quantum systems are synchronized with coherent
vacuum fluctuations, the total vacuum energy scales as
$E \propto N^2$ (analogous to superradiance~\cite{Dicke1954}).
For $N = 10^{15}$ synchronized resonators:
$E \sim 10^{-22} \times 10^{30} = 10^{8}\,\mathrm{J}$,
approaching the Planck energy.

\textit{Open problem I:} Under what conditions do vacuum energy
modifications exhibit $N^2$ scaling rather than $N$ scaling?
What is the decoherence rate for vacuum-energy coherence?

\textbf{Mechanism II: Topological protection.}
If the sign flip of vacuum energy is a topological invariant
(protected by the K-theory classification of phases), it may
be robust against thermal fluctuations and decoherence.
The relevant K-group change is
$K_1(\mathbb{Z}) = \mathbb{Z}/2 \to K_1(\Zp) = \mathbb{Z}/2 \times \mathbb{Z}$,
a topological phase transition.

\textit{Open problem II:} Does the K-theoretic phase transition
$K_1(\mathbb{Z}) \to K_1(\Zp)$ correspond to a
physically realizable topological phase transition?
If so, are the associated edge states macroscopically
observable?

\textbf{Mechanism III: Critical amplification.}
Near the BC phase transition at $\beta = 1$, fluctuations
diverge: $\langle (\Delta E)^2 \rangle \propto |\beta - 1|^{-\gamma}$
for some critical exponent $\gamma$. If the vacuum energy
modification is coupled to these critical fluctuations, it
could be amplified near the transition point.

\textit{Open problem III:} What are the critical exponents
of the BC system at $\beta = 1$? Does the prime muting
operation couple to the critical mode?

\textbf{Mechanism IV: Holographic amplification.}
In AdS/CFT, boundary perturbations can produce
$O(1/G_N) \sim O(N^2)$ effects in the bulk, where $N$
is the rank of the gauge group and $G_N$ is Newton's constant.
If the BC system is a boundary theory for an arithmetic bulk,
prime muting on the boundary could produce amplified effects
in the bulk geometry.

\textit{Open problem IV:} Is there an arithmetic analog of
AdS/CFT where the BC system serves as boundary theory?
What is the ``central charge'' of such a correspondence?

\textbf{Mechanism V: Bootstrap.}
If a Planck-scale warp bubble can be created
($E \sim 10^9\,\mathrm{J}$, potentially reachable
via Mechanisms I--IV), it might serve as a seed for
a self-reinforcing process: the bubble itself modifies
spacetime in a way that reduces the energy cost of
further expansion.

\textit{Open problem V:} Is a Planck-scale warp bubble
stable? Can it serve as a seed for controlled expansion?
What is the energy cost as a function of bubble radius?

\subsection{Qualitative vs.\ quantitative resolution}

We note that the scaling problem may admit a
\emph{qualitative} rather than quantitative resolution.
In the companion paper on sheaf-theoretic
reformulation~\cite{paper_B}, we argue that the WEC
may be interpretable as a cohomological obstruction
on $\Spec(\mathbb{Z})$ whose vanishing under localization
$\mathbb{Z} \to \Zp$ is a discrete (on/off) transition
rather than a continuous energy threshold. If this
interpretation is correct, the scaling problem dissolves:
one does not need to ``accumulate'' negative energy, but
rather to ``switch off'' the obstruction.

This remains a conjecture (see~\cite{paper_B}, Conjecture~3).

% ============================================================================
\section{The Role of Experimental Verification}
\label{sec:experiment}
% ============================================================================

The experimental proposal in~\cite{paper_A} directly tests
Steps~4--6 of the chain (Table~\ref{tab:chain}).
We summarize the relationship between experimental outcomes
and the chain:

\begin{description}
  \item[Level 1 success] ($\Delta E \neq 0$ measured)
    $\Rightarrow$ Steps~4--5 confirmed.
    Vacuum energy depends on mode arithmetic.

  \item[Level 2 success] ($\Delta E$ matches $\zeta$-regularized value)
    $\Rightarrow$ Step~6 confirmed.
    Zeta regularization correctly predicts the physical vacuum energy
    under mode selection.

  \item[Level 3 success] (multiplicative structure verified for $p=2,3$)
    $\Rightarrow$ Step~5 confirmed as a physical law.
    The Euler product is not merely a mathematical identity but
    governs the quantum vacuum.

  \item[Level 4 success] (controllable sign-flipping demonstrated)
    $\Rightarrow$ Step~7 confirmed at laboratory scale.
    WEC violation via arithmetic control is physically possible.
\end{description}

Steps~2 and~9 (the Argued steps) are not directly testable by
the proposed experiment. They require either:
(a) independent confirmation of the NCG Standard Model
(e.g., prediction of a new particle or coupling constant), or
(b) detection of a gravitational signature of vacuum energy
modification (requiring sensitivity beyond current technology).

Step~10 (Scaling) is not testable by any currently proposed
experiment. It requires theoretical advances in one or more of
Mechanisms~I--V above.

% ============================================================================
\section{Relation to Other Warp Drive Proposals}
\label{sec:related}
% ============================================================================

\textbf{Lentz solitons~\cite{Lentz2021}.}
Lentz showed that subluminal warp-like solitons with purely
positive energy exist in classical GR. These solutions do not
require WEC violation. Our approach complements Lentz's:
if superluminal travel is the goal, WEC violation is necessary
(as shown by Bobrick-Martire~\cite{BobrickMartire2021}),
and our mechanism provides a specific algebraic route to it.
For subluminal applications, Lentz solitons may be achievable
without arithmetic vacuum engineering.

\textbf{Bobrick-Martire classification~\cite{BobrickMartire2021}.}
Bobrick and Martire classified warp drives as shell geometries
around flat Minkowski interiors. Their key result: subluminal
positive-energy warp is possible; superluminal requires WEC
violation. Our chain addresses the superluminal case by providing
a WEC-violating mechanism (Steps~4--7).

\textbf{White (EagleWorks).}
White~\cite{White2013} proposed interferometric detection of
microscopic warp-field perturbations using electromagnetic
cavities. Our SQUID-based experiment (Section~VI of~\cite{paper_A})
is conceptually related but tests a more specific hypothesis:
the arithmetic (Euler product) structure of vacuum energy.

% ============================================================================
\section{Conclusion and Roadmap}
\label{sec:conclusion}
% ============================================================================

The theoretical chain from the BC system to warp drive consists
of ten steps (Table~\ref{tab:chain}). Seven are established,
two are argued, and one is open:

\begin{enumerate}
  \item \textbf{Established} (Steps 1, 3--8): The mathematical
    framework is rigorous. The sign-flip theorem is an algebraic
    identity. The physical interpretation via Casimir effect is
    standard.

  \item \textbf{Argued} (Steps 2, 9): The connection between the
    BC system and physical spacetime relies on the NCG program,
    which successfully reproduces the Standard Model but lacks
    independent experimental confirmation. These steps can be
    formulated as precise conjectures
    (Conjecture~\ref{conj:kernel} and Bridge~A/B above).

  \item \textbf{Open} (Step 10): The scaling problem is the most
    critical gap. Five candidate mechanisms are identified, each
    with a well-defined open problem.
\end{enumerate}

The recommended research roadmap is:

\textbf{Phase 1 (1--2 years):}
Execute the SQUID experiment~\cite{paper_A} to test Steps~4--6.

\textbf{Phase 2 (2--5 years):}
Develop the sheaf-theoretic framework~\cite{paper_B} to
address the qualitative resolution of the scaling problem.
Investigate topological phase transitions (Mechanism~II) in
condensed matter systems.

\textbf{Phase 3 (5--10 years):}
Pursue coherent amplification (Mechanism~I) using arrays of
synchronized quantum circuits. Investigate the critical behavior
of the BC system (Mechanism~III).

\textbf{Phase 4 (10+ years):}
If Phase~1 confirms the arithmetic structure of vacuum energy,
and Phase~2--3 identify a viable amplification mechanism,
pursue macroscopic negative energy generation.

The arithmetic vacuum engineering program is ambitious but
scientifically grounded: each step is either proven, precisely
conjectured, or formulated as a well-defined open problem.
Even partial success --- confirming the Euler product structure
of vacuum energy --- would be a significant contribution to
fundamental physics.

\begin{acknowledgments}
This work was conducted as independent research by the
Wright Brothers collaboration.
\end{acknowledgments}

\begin{thebibliography}{30}

\bibitem{paper_A}
Wright Brothers,
``Arithmetic vacuum engineering: Prime channel muting via
Bost-Connes phase transitions and its implications for
exotic matter generation,''
companion paper (2026).

\bibitem{paper_B}
Wright Brothers,
``Sheaf-theoretic reformulation of vacuum energy conditions
on $\Spec(\mathbb{Z})$,''
companion paper (2026).

\bibitem{Alcubierre1994}
M.~Alcubierre,
Class. Quantum Grav. \textbf{11}, L73 (1994).

\bibitem{BostConnes1995}
J.-B.~Bost and A.~Connes,
Selecta Math. (N.S.) \textbf{1}, 411 (1995).

\bibitem{ConnesMarcolli2006}
A.~Connes and M.~Marcolli,
\emph{Noncommutative Geometry, Quantum Fields and Motives},
AMS (2008); see Theorem~4.1 and Chapter~3.

\bibitem{Connes1996}
A.~Connes,
``Gravity coupled with matter and the foundation of
non-commutative geometry,''
Commun. Math. Phys. \textbf{182}, 155 (1996).

\bibitem{CCM2007}
A.~H.~Chamseddine, A.~Connes, and M.~Marcolli,
``Gravity and the Standard Model with neutrino mixing,''
Adv. Theor. Math. Phys. \textbf{11}, 991 (2007).

\bibitem{Casimir1948}
H.~B.~G.~Casimir,
Proc. K. Ned. Akad. Wet. \textbf{51}, 793 (1948).

\bibitem{Bordag2009}
M.~Bordag et al.,
\emph{Advances in the Casimir Effect},
Oxford University Press (2009).

\bibitem{Maldacena1998}
J.~Maldacena,
``The large $N$ limit of superconformal field theories
and supergravity,''
Adv. Theor. Math. Phys. \textbf{2}, 231 (1998).

\bibitem{Dicke1954}
R.~H.~Dicke,
``Coherence in spontaneous radiation processes,''
Phys. Rev. \textbf{93}, 99 (1954).

\bibitem{Lentz2021}
E.~W.~Lentz,
Class. Quantum Grav. \textbf{38}, 075015 (2021).

\bibitem{BobrickMartire2021}
A.~Bobrick and G.~Martire,
Class. Quantum Grav. \textbf{38}, 105009 (2021).

\bibitem{White2013}
H.~White,
``Warp field mechanics 102: Energy optimization,''
NASA Technical Report JSC-CN-28555 (2013).

\bibitem{Lamoreaux1997}
S.~K.~Lamoreaux,
``Demonstration of the Casimir force in the 0.6 to $6\,\mu$m range,''
Phys. Rev. Lett. \textbf{78}, 5 (1997).

\end{thebibliography}

\end{document}
